% Options for packages loaded elsewhere
\PassOptionsToPackage{unicode}{hyperref}
\PassOptionsToPackage{hyphens}{url}
\documentclass[
  ignorenonframetext,
]{beamer}
\newif\ifbibliography
\usepackage{pgfpages}
\setbeamertemplate{caption}[numbered]
\setbeamertemplate{caption label separator}{: }
\setbeamercolor{caption name}{fg=normal text.fg}
\beamertemplatenavigationsymbolsempty
% remove section numbering
\setbeamertemplate{part page}{
  \centering
  \begin{beamercolorbox}[sep=16pt,center]{part title}
    \usebeamerfont{part title}\insertpart\par
  \end{beamercolorbox}
}
\setbeamertemplate{section page}{
  \centering
  \begin{beamercolorbox}[sep=12pt,center]{section title}
    \usebeamerfont{section title}\insertsection\par
  \end{beamercolorbox}
}
\setbeamertemplate{subsection page}{
  \centering
  \begin{beamercolorbox}[sep=8pt,center]{subsection title}
    \usebeamerfont{subsection title}\insertsubsection\par
  \end{beamercolorbox}
}
% Prevent slide breaks in the middle of a paragraph
\widowpenalties 1 10000
\raggedbottom
\AtBeginPart{
  \frame{\partpage}
}
\AtBeginSection{
  \ifbibliography
  \else
    \frame{\sectionpage}
  \fi
}
\AtBeginSubsection{
  \frame{\subsectionpage}
}
\usepackage{iftex}
\ifPDFTeX
  \usepackage[T1]{fontenc}
  \usepackage[utf8]{inputenc}
  \usepackage{textcomp} % provide euro and other symbols
\else % if luatex or xetex
  \usepackage{unicode-math} % this also loads fontspec
  \defaultfontfeatures{Scale=MatchLowercase}
  \defaultfontfeatures[\rmfamily]{Ligatures=TeX,Scale=1}
\fi
\usepackage{lmodern}
\ifPDFTeX\else
  % xetex/luatex font selection
\fi
% Use upquote if available, for straight quotes in verbatim environments
\IfFileExists{upquote.sty}{\usepackage{upquote}}{}
\IfFileExists{microtype.sty}{% use microtype if available
  \usepackage[]{microtype}
  \UseMicrotypeSet[protrusion]{basicmath} % disable protrusion for tt fonts
}{}
\makeatletter
\@ifundefined{KOMAClassName}{% if non-KOMA class
  \IfFileExists{parskip.sty}{%
    \usepackage{parskip}
  }{% else
    \setlength{\parindent}{0pt}
    \setlength{\parskip}{6pt plus 2pt minus 1pt}}
}{% if KOMA class
  \KOMAoptions{parskip=half}}
\makeatother
\usepackage{color}
\usepackage{fancyvrb}
\newcommand{\VerbBar}{|}
\newcommand{\VERB}{\Verb[commandchars=\\\{\}]}
\DefineVerbatimEnvironment{Highlighting}{Verbatim}{commandchars=\\\{\}}
% Add ',fontsize=\small' for more characters per line
\usepackage{framed}
\definecolor{shadecolor}{RGB}{248,248,248}
\newenvironment{Shaded}{\begin{snugshade}}{\end{snugshade}}
\newcommand{\AlertTok}[1]{\textcolor[rgb]{0.94,0.16,0.16}{#1}}
\newcommand{\AnnotationTok}[1]{\textcolor[rgb]{0.56,0.35,0.01}{\textbf{\textit{#1}}}}
\newcommand{\AttributeTok}[1]{\textcolor[rgb]{0.13,0.29,0.53}{#1}}
\newcommand{\BaseNTok}[1]{\textcolor[rgb]{0.00,0.00,0.81}{#1}}
\newcommand{\BuiltInTok}[1]{#1}
\newcommand{\CharTok}[1]{\textcolor[rgb]{0.31,0.60,0.02}{#1}}
\newcommand{\CommentTok}[1]{\textcolor[rgb]{0.56,0.35,0.01}{\textit{#1}}}
\newcommand{\CommentVarTok}[1]{\textcolor[rgb]{0.56,0.35,0.01}{\textbf{\textit{#1}}}}
\newcommand{\ConstantTok}[1]{\textcolor[rgb]{0.56,0.35,0.01}{#1}}
\newcommand{\ControlFlowTok}[1]{\textcolor[rgb]{0.13,0.29,0.53}{\textbf{#1}}}
\newcommand{\DataTypeTok}[1]{\textcolor[rgb]{0.13,0.29,0.53}{#1}}
\newcommand{\DecValTok}[1]{\textcolor[rgb]{0.00,0.00,0.81}{#1}}
\newcommand{\DocumentationTok}[1]{\textcolor[rgb]{0.56,0.35,0.01}{\textbf{\textit{#1}}}}
\newcommand{\ErrorTok}[1]{\textcolor[rgb]{0.64,0.00,0.00}{\textbf{#1}}}
\newcommand{\ExtensionTok}[1]{#1}
\newcommand{\FloatTok}[1]{\textcolor[rgb]{0.00,0.00,0.81}{#1}}
\newcommand{\FunctionTok}[1]{\textcolor[rgb]{0.13,0.29,0.53}{\textbf{#1}}}
\newcommand{\ImportTok}[1]{#1}
\newcommand{\InformationTok}[1]{\textcolor[rgb]{0.56,0.35,0.01}{\textbf{\textit{#1}}}}
\newcommand{\KeywordTok}[1]{\textcolor[rgb]{0.13,0.29,0.53}{\textbf{#1}}}
\newcommand{\NormalTok}[1]{#1}
\newcommand{\OperatorTok}[1]{\textcolor[rgb]{0.81,0.36,0.00}{\textbf{#1}}}
\newcommand{\OtherTok}[1]{\textcolor[rgb]{0.56,0.35,0.01}{#1}}
\newcommand{\PreprocessorTok}[1]{\textcolor[rgb]{0.56,0.35,0.01}{\textit{#1}}}
\newcommand{\RegionMarkerTok}[1]{#1}
\newcommand{\SpecialCharTok}[1]{\textcolor[rgb]{0.81,0.36,0.00}{\textbf{#1}}}
\newcommand{\SpecialStringTok}[1]{\textcolor[rgb]{0.31,0.60,0.02}{#1}}
\newcommand{\StringTok}[1]{\textcolor[rgb]{0.31,0.60,0.02}{#1}}
\newcommand{\VariableTok}[1]{\textcolor[rgb]{0.00,0.00,0.00}{#1}}
\newcommand{\VerbatimStringTok}[1]{\textcolor[rgb]{0.31,0.60,0.02}{#1}}
\newcommand{\WarningTok}[1]{\textcolor[rgb]{0.56,0.35,0.01}{\textbf{\textit{#1}}}}
\usepackage{graphicx}
\makeatletter
\newsavebox\pandoc@box
\newcommand*\pandocbounded[1]{% scales image to fit in text height/width
  \sbox\pandoc@box{#1}%
  \Gscale@div\@tempa{\textheight}{\dimexpr\ht\pandoc@box+\dp\pandoc@box\relax}%
  \Gscale@div\@tempb{\linewidth}{\wd\pandoc@box}%
  \ifdim\@tempb\p@<\@tempa\p@\let\@tempa\@tempb\fi% select the smaller of both
  \ifdim\@tempa\p@<\p@\scalebox{\@tempa}{\usebox\pandoc@box}%
  \else\usebox{\pandoc@box}%
  \fi%
}
% Set default figure placement to htbp
\def\fps@figure{htbp}
\makeatother
\setlength{\emergencystretch}{3em} % prevent overfull lines
\providecommand{\tightlist}{%
  \setlength{\itemsep}{0pt}\setlength{\parskip}{0pt}}
\usepackage{bookmark}
\IfFileExists{xurl.sty}{\usepackage{xurl}}{} % add URL line breaks if available
\urlstyle{same}
\hypersetup{
  pdftitle={Regression 1},
  pdfauthor={Aliaksandr Hubin},
  hidelinks,
  pdfcreator={LaTeX via pandoc}}

\title{Regression 1}
\author{Aliaksandr Hubin}
\date{February 10 2025}

\begin{document}
\frame{\titlepage}

\begin{frame}[fragile]{The body data}
\phantomsection\label{the-body-data}
A motivating example - The body data

\begin{Shaded}
\begin{Highlighting}[]
\FunctionTok{library}\NormalTok{(BIAS.data)}
\FunctionTok{head}\NormalTok{(bodydata,}\DecValTok{10}\NormalTok{)}
\end{Highlighting}
\end{Shaded}

\begin{verbatim}
##    Weight Height Age Circumference
## 1    65.6  174.0  21          71.5
## 3    80.7  193.5  28          83.2
## 4    72.6  186.5  23          77.8
## 5    78.8  187.2  22          80.0
## 6    74.8  181.5  21          82.5
## 7    86.4  184.0  26          82.0
## 8    78.4  184.5  27          76.8
## 9    62.0  175.0  23          68.5
## 10   81.6  184.0  21          77.5
## 11   76.6  180.0  23          81.9
\end{verbatim}
\end{frame}

\begin{frame}{Scatterplots}
\phantomsection\label{scatterplots}
\pandocbounded{\includegraphics[keepaspectratio]{Lecture_2_files/figure-beamer/unnamed-chunk-2-1.pdf}}
\end{frame}

\begin{frame}{The Weight versus Height}
\phantomsection\label{the-weight-versus-height}
Is there a linear relationship between Weight and Height?

\pandocbounded{\includegraphics[keepaspectratio]{Lecture_2_files/figure-beamer/unnamed-chunk-3-1.pdf}}

If so, what is the ``best'' fitting line?
\end{frame}

\begin{frame}{A graphical presentation of a simple linear model}
\phantomsection\label{a-graphical-presentation-of-a-simple-linear-model}
Let's use a web applet to discuss the ``best'' fitting line:
\url{http://solve.shinyapps.io/LeastSquaresApp}
\end{frame}

\begin{frame}{The main idea}
\phantomsection\label{the-main-idea}
A continuous response variable, \(y\), is assumed to be linearly
dependent on one or several predictor variables,
\(x_1, x_2, \ldots, x_p\)

A simple regression model (one predictor)

\(y = \beta_0 + \beta_1 x + \epsilon, \epsilon \sim N(0, \sigma^2)\)

A multiple regression model (multiple predictors)

\(y = \beta_0 + \beta_1 x_1 + \beta_2 x_2 + \cdots+ \beta_p x_p +\epsilon, \epsilon \sim N(0, \sigma^2)\)
\end{frame}

\begin{frame}{Estimation of regression coefficients}
\phantomsection\label{estimation-of-regression-coefficients}
Unknown regression coefficients of a simple regression model are
\(\beta_0\) and \(\beta_1\)

Estimation of parameters is equivalent to finding the ``best fitting''
model.

Different estimation methods \(\Leftrightarrow\) Different definitions
of ``best''.

\begin{itemize}
\tightlist
\item
  Least squares estimator : The estimates minimize
  \(SSE(\beta_0,\beta_1) = \sum (y_i - ({\beta}_0 + {\beta}_1 x_i))^2\rightarrow\min\)
\item
  Equivalent to the maximum likelihood estimator: The estimates make the
  observed data ``most likely''
\item
  \[
  f(y_1,...,y_n)\rightarrow\max
  \]
\end{itemize}
\end{frame}

\begin{frame}{Maximum likelihood formulation}
\phantomsection\label{maximum-likelihood-formulation}
Recall normal
pdf\(f(y) = \frac{1}{\sigma\sqrt{2\pi}} \exp\left\{-\frac{1}{2}\left(\frac{y-\mu}{\sigma}\right)^2\right\}\)
\end{frame}

\begin{frame}{Derivations for maxima}
\phantomsection\label{derivations-for-maxima}
\end{frame}

\begin{frame}{Estimation of unexplained variance}
\phantomsection\label{estimation-of-unexplained-variance}
As soon as we have defined the ``best'' fitting model, we have also
defined all distances between the observations and the estimated model.

These distances are the residuals:
\[d_i = y_i - (\widehat{\beta}_0 + \widehat{\beta}_1 x_i)\] for
\(i=1,\ldots,n\)

An unbiased estimator for the error variance \(\sigma^2\) depends on
these distances through
\[ \widehat{\sigma}^2 = \sum d_i^2/(n-p-1) = \frac{SSE}{(n-p-1)} = MSE\]

where \(p\) is the number of predictors.
\end{frame}

\begin{frame}{Example: Regression in R (Model 1)}
\phantomsection\label{example-regression-in-r-model-1}
We use samples 1 to 20 of the bodydata as example data (training set)

\[ y_i = \beta_0 +\beta_1 x_{1i} + \epsilon_i \] where \(y\) is Weight,
\(x_1\) is Height and \(i=1,\ldots, 20\).

\pandocbounded{\includegraphics[keepaspectratio]{Lecture_2_files/figure-beamer/unnamed-chunk-4-1.pdf}}
\end{frame}

\begin{frame}[fragile]{Example of model fit in R Studio}
\phantomsection\label{example-of-model-fit-in-r-studio}
\begin{Shaded}
\begin{Highlighting}[]
\NormalTok{LinearModel}\FloatTok{.1} \OtherTok{\textless{}{-}} \FunctionTok{lm}\NormalTok{(Weight }\SpecialCharTok{\textasciitilde{}}\NormalTok{ Height, }\AttributeTok{data=}\NormalTok{bodydata, }\AttributeTok{subset=}\DecValTok{1}\SpecialCharTok{:}\DecValTok{20}\NormalTok{)}
\FunctionTok{summary}\NormalTok{(LinearModel}\FloatTok{.1}\NormalTok{)}
\end{Highlighting}
\end{Shaded}
\end{frame}

\begin{frame}[fragile]{Output}
\phantomsection\label{output}
\begin{verbatim}
## 
## Call:
## lm(formula = Weight ~ Height, data = bodydata, subset = 1:20)
## 
## Residuals:
##      Min       1Q   Median       3Q      Max 
## -11.5160  -2.5964   0.0261   2.9141  11.7454 
## 
## Coefficients:
##             Estimate Std. Error t value Pr(>|t|)   
## (Intercept) -38.2310    38.3605  -0.997  0.33216   
## Height        0.6386     0.2123   3.007  0.00757 **
## ---
## Signif. codes:  0 '***' 0.001 '**' 0.01 '*' 0.05 '.' 0.1 ' ' 1
## 
## Residual standard error: 5.839 on 18 degrees of freedom
## Multiple R-squared:  0.3344, Adjusted R-squared:  0.2974 
## F-statistic: 9.043 on 1 and 18 DF,  p-value: 0.007566
\end{verbatim}
\end{frame}

\begin{frame}[fragile]{The fitted line plot}
\phantomsection\label{the-fitted-line-plot}
\begin{Shaded}
\begin{Highlighting}[]
\FunctionTok{plot}\NormalTok{(Weight}\SpecialCharTok{\textasciitilde{}}\NormalTok{Height, }\AttributeTok{data=}\NormalTok{bodydata, }\AttributeTok{subset=}\DecValTok{1}\SpecialCharTok{:}\DecValTok{20}\NormalTok{, }\AttributeTok{pch=}\DecValTok{20}\NormalTok{, }\AttributeTok{cex=}\DecValTok{2}\NormalTok{)}
\FunctionTok{abline}\NormalTok{(LinearModel}\FloatTok{.1}\NormalTok{, }\AttributeTok{lwd=}\DecValTok{2}\NormalTok{, }\AttributeTok{col=}\StringTok{"red"}\NormalTok{)}
\end{Highlighting}
\end{Shaded}

\pandocbounded{\includegraphics[keepaspectratio]{Lecture_2_files/figure-beamer/unnamed-chunk-7-1.pdf}}
\end{frame}

\begin{frame}{Weight versus Circumference (Model 2)}
\phantomsection\label{weight-versus-circumference-model-2}
Now let:

\[ y_i = \beta_0 +\beta_1 x_{2i} + \epsilon_i \]

where \(y\) still is Weight, \(x_2\) is Circumference and
\(i=1,\ldots, 20\).

\pandocbounded{\includegraphics[keepaspectratio]{Lecture_2_files/figure-beamer/unnamed-chunk-8-1.pdf}}
\end{frame}

\begin{frame}[fragile]{R Commander output}
\phantomsection\label{r-commander-output}
\begin{verbatim}
## 
## Call:
## lm(formula = Weight ~ Circumference, data = bodydata, subset = 1:20)
## 
## Residuals:
##    Min     1Q Median     3Q    Max 
## -8.937 -3.540  0.038  2.956  8.659 
## 
## Coefficients:
##               Estimate Std. Error t value Pr(>|t|)    
## (Intercept)     3.6737    17.1036   0.215 0.832343    
## Circumference   0.9137     0.2125   4.300 0.000431 ***
## ---
## Signif. codes:  0 '***' 0.001 '**' 0.01 '*' 0.05 '.' 0.1 ' ' 1
## 
## Residual standard error: 5.027 on 18 degrees of freedom
## Multiple R-squared:  0.5067, Adjusted R-squared:  0.4793 
## F-statistic: 18.49 on 1 and 18 DF,  p-value: 0.0004311
\end{verbatim}

How can we claim that this model gives a better fit than model1?
\end{frame}

\begin{frame}{A measure of goodness of fit, the \(R^2\)}
\phantomsection\label{a-measure-of-goodness-of-fit-the-r2}
The unconditional (sample) variance of Weight is:
\[\sum (y_i - \bar{y})^2 = {SST}\] Hence, \(SST\) (in some literature
\(SSY\)) is a mesure of total variability in y, the response

Further, \(SSE\) is a measure of unexplained variability in y

So, a measure of explained variability is
\[R^2 = \frac{SST-SSE}{SST} = 1-\frac{SSE}{SST}\]
\end{frame}

\begin{frame}{Multiple linear regression (Model 3)}
\phantomsection\label{multiple-linear-regression-model-3}
\[y_i = \beta_0 + \beta_1 x_{1i} + \beta_2x_{2i} + \epsilon_i\]

where \(y\) is Weight, \(x_1\) is Height \(x_2\) is Circumference and
\(i=1,\ldots, 20\).
\end{frame}

\begin{frame}[fragile]{R Studio output (Model 3)}
\phantomsection\label{r-studio-output-model-3}
\begin{verbatim}
## 
## Call:
## lm(formula = Weight ~ Height + Circumference, data = bodydata, 
##     subset = 1:20)
## 
## Residuals:
##     Min      1Q  Median      3Q     Max 
## -5.3189 -3.5363 -0.7817  2.8028  6.3974 
## 
## Coefficients:
##               Estimate Std. Error t value Pr(>|t|)    
## (Intercept)   -70.8202    28.4518  -2.489 0.023464 *  
## Height          0.4731     0.1566   3.021 0.007698 ** 
## Circumference   0.7777     0.1820   4.273 0.000515 ***
## ---
## Signif. codes:  0 '***' 0.001 '**' 0.01 '*' 0.05 '.' 0.1 ' ' 1
## 
## Residual standard error: 4.172 on 17 degrees of freedom
## Multiple R-squared:  0.679,  Adjusted R-squared:  0.6413 
## F-statistic: 17.98 on 2 and 17 DF,  p-value: 6.38e-05
\end{verbatim}
\end{frame}

\begin{frame}[fragile]{Three predictors\ldots(Model 4)}
\phantomsection\label{three-predictorsmodel-4}
\begin{verbatim}
## 
## Call:
## lm(formula = Weight ~ Height + Circumference + Age, data = bodydata, 
##     subset = 1:20)
## 
## Residuals:
##     Min      1Q  Median      3Q     Max 
## -5.4088 -3.3900 -0.9123  3.0403  6.9827 
## 
## Coefficients:
##               Estimate Std. Error t value Pr(>|t|)    
## (Intercept)   -66.5785    30.0098  -2.219 0.041330 *  
## Height          0.4768     0.1600   2.981 0.008831 ** 
## Circumference   0.7769     0.1858   4.181 0.000706 ***
## Age            -0.2094     0.3731  -0.561 0.582352    
## ---
## Signif. codes:  0 '***' 0.001 '**' 0.01 '*' 0.05 '.' 0.1 ' ' 1
## 
## Residual standard error: 4.259 on 16 degrees of freedom
## Multiple R-squared:  0.6852, Adjusted R-squared:  0.6262 
## F-statistic: 11.61 on 3 and 16 DF,  p-value: 0.0002729
\end{verbatim}
\end{frame}

\begin{frame}{\(R^2\)-adjusted}
\phantomsection\label{r2-adjusted}
\[R^2 = 1 - \frac{SSE}{SST} \text{, but the Problem: $R^2$ always increases when adding predictors}\]

\[R^2_{\text{adj}} = 1 - \frac{\widehat\sigma^2}{\widehat{Var}(y)} = 1 - \frac{SSE/(n-p-1)}{SST/(n-1)}\]
The Solution: \(R^2_{\text{adjusted}}\) penalizes adding parameters

\includegraphics[width=400px]{Lecture_2_files/figure-beamer/unnamed-chunk-12-1}
\end{frame}

\begin{frame}{Hypothesis Test for Regression Coefficients}
\phantomsection\label{hypothesis-test-for-regression-coefficients}
We may test the \textbf{conditional significance} of any predictor using
a t-test

\textbf{Hypotheses:}

\begin{itemize}
\item
  \(H_0: \beta_j = 0\) (predictor has \textbf{no effect} on response)
\item
  \(H_1: \beta_j \neq 0\) (predictor \textbf{does have an effect})
\end{itemize}

\textbf{Test statistic:}
\[T_j = \frac{\widehat{\beta}_j-0}{SE(\widehat{\beta}_j)} \sim t_{n-p-1}\]
\end{frame}

\begin{frame}{Visual explanation}
\phantomsection\label{visual-explanation}
\end{frame}

\begin{frame}{Decision rule}
\phantomsection\label{decision-rule}
The \textbf{p-value} is the probability of observing a test statistic as
extreme or more extreme than the observed value, \textbf{assuming}
\(H_0\) is true

\[\text{p-value} = P(|T| \geq |T_{\text{obs}}|) = 2 \times P(T > |T_{\text{obs}}|)\]

\begin{itemize}
\tightlist
\item
  \textbf{p-value is small} (e.g., 0.002): It's \textbf{very unlikely}
  to observe such an extreme
\item
  T if \(H_0\) were true → Strong evidence against \(H_0\) →
  \textbf{Reject} \(H_0\)
\end{itemize}

\textbf{At significance level of 0.05:}

\begin{itemize}
\tightlist
\item
  Reject \(H_0\) if \(|T_j| > \tau_{\alpha/2, n-p-1}\) (i.e., p-value
  \(< 0.05\))
\item
  Fail to reject \(H_0\) if \(|T_j| \leq \tau_{\alpha/2, n-p-1}\) (i.e.,
  p-value \(\geq 0.05\))
\end{itemize}
\end{frame}

\begin{frame}[fragile]{Example: Model 4 Significance tests}
\phantomsection\label{example-model-4-significance-tests}
\begin{verbatim}
##                  Estimate Std. Error    t value     Pr(>|t|)
## (Intercept)   -66.5785315 30.0098013 -2.2185596 0.0413297267
## Height          0.4768387  0.1599838  2.9805444 0.0088305547
## Circumference   0.7768555  0.1858173  4.1807501 0.0007063713
## Age            -0.2094104  0.3730630 -0.5613273 0.5823518801
\end{verbatim}

\textbf{Key insight:}

\begin{itemize}
\item
  Age is not significant (p = 0.582 \textgreater\textgreater{} 0.05)
\item
  Explains why adding Age slightly reduces \(R^2_{\text{adj}}\) - Model
  3 (without Age) is \textbf{better choice} than Model 4
\item
  Other predictors are significant
\end{itemize}
\end{frame}

\begin{frame}[fragile]{Confidence interval for regression coefficients}
\phantomsection\label{confidence-interval-for-regression-coefficients}
A \((1-\alpha)\cdot 100 \%\) for \(\beta_j\) is mathematically given by:
\[ \left[ \widehat{\beta}_j\pm \tau_{\alpha/2, (n-p-1)} SE(\widehat{\beta}_j) \right]\]

In RStudio we can get these (ex. for model 4) by

\begin{Shaded}
\begin{Highlighting}[]
\FunctionTok{confint}\NormalTok{(LinearModel}\FloatTok{.4}\NormalTok{, }\AttributeTok{level=}\FloatTok{0.95}\NormalTok{)}
\end{Highlighting}
\end{Shaded}

\begin{verbatim}
##                      2.5 %     97.5 %
## (Intercept)   -130.1964683 -2.9605946
## Height           0.1376883  0.8159891
## Circumference    0.3829405  1.1707706
## Age             -1.0002686  0.5814477
\end{verbatim}
\end{frame}

\begin{frame}[fragile]{Prediction}
\phantomsection\label{prediction}
A fitted model can be used for prediction of new observations

According to model 4, what is the predicted Weight of some person with
these values?

\begin{verbatim}
##   Height Circumference Age
## 1    171            76  30
\end{verbatim}

\[
\begin{split}
\widehat{y} &= -66.579 + 0.4768\cdot171 + 0.7769\cdot76 - 0.2094\cdot30 \\
& = 67.72
\end{split}
\]

\begin{Shaded}
\begin{Highlighting}[]
\NormalTok{newobs }\OtherTok{\textless{}{-}} \FunctionTok{data.frame}\NormalTok{(}\StringTok{\textquotesingle{}Height\textquotesingle{}}\OtherTok{=}\DecValTok{171}\NormalTok{, }\StringTok{\textquotesingle{}Circumference\textquotesingle{}}\OtherTok{=}\DecValTok{76}\NormalTok{, }\StringTok{\textquotesingle{}Age\textquotesingle{}}\OtherTok{=}\DecValTok{30}\NormalTok{)}
\FunctionTok{predict.lm}\NormalTok{(LinearModel}\FloatTok{.4}\NormalTok{, }\AttributeTok{newdata=}\NormalTok{newobs)}
\end{Highlighting}
\end{Shaded}

\begin{verbatim}
##       1 
## 67.7196
\end{verbatim}
\end{frame}

\begin{frame}[fragile]{Other intervals in Prediction output}
\phantomsection\label{other-intervals-in-prediction-output}
Confidence interval for \(E(y|x_1, \ldots, x_p)\)

\begin{Shaded}
\begin{Highlighting}[]
\FunctionTok{predict.lm}\NormalTok{(LinearModel}\FloatTok{.4}\NormalTok{, }\AttributeTok{newdata=}\NormalTok{newobs, }
           \AttributeTok{level=}\FloatTok{0.95}\NormalTok{, }\AttributeTok{interval=}\FunctionTok{c}\NormalTok{(}\StringTok{"confidence"}\NormalTok{))}
\end{Highlighting}
\end{Shaded}

\begin{verbatim}
##       fit      lwr      upr
## 1 67.7196 60.98252 74.45668
\end{verbatim}

Prediction interval for \(y|x_1, \ldots, x_p\)

\begin{Shaded}
\begin{Highlighting}[]
\FunctionTok{predict.lm}\NormalTok{(LinearModel}\FloatTok{.4}\NormalTok{, }\AttributeTok{newdata=}\NormalTok{newobs, }
           \AttributeTok{level=}\FloatTok{0.95}\NormalTok{, }\AttributeTok{interval=}\FunctionTok{c}\NormalTok{(}\StringTok{"prediction"}\NormalTok{))}
\end{Highlighting}
\end{Shaded}

\begin{verbatim}
##       fit      lwr      upr
## 1 67.7196 56.45463 78.98456
\end{verbatim}
\end{frame}

\begin{frame}[fragile]{Compare real and predicted values}
\phantomsection\label{compare-real-and-predicted-values}
Let's predict the 387 extra observations and compare with actual
observations!

\begin{verbatim}
##         fit      lwr      upr
## 25 73.94436 64.06824 83.82047
## 26 82.80148 72.72603 92.87693
## 27 81.03059 71.44393 90.61725
## 28 71.46852 61.81893 81.11811
## 29 65.93065 55.27452 76.58677
## 31 76.31124 66.90951 85.71296
## 32 69.35871 59.61763 79.09979
## 33 75.73802 65.67258 85.80346
## 34 82.52507 72.84459 92.20555
## 35 67.71321 57.86212 77.56431
## 37 73.86325 64.29154 83.43496
## 39 74.81258 65.46899 84.15617
\end{verbatim}
\end{frame}

\begin{frame}{Predicted vs observed - plot}
\phantomsection\label{predicted-vs-observed---plot}
\includegraphics[width=700px,height=400]{Lecture_2_files/figure-beamer/unnamed-chunk-20-1}
\end{frame}

\begin{frame}[fragile]{Checking the prediction intervals}
\phantomsection\label{checking-the-prediction-intervals}
How many observations fell within the prediction intervals?

\begin{Shaded}
\begin{Highlighting}[]
\NormalTok{largerthanlower }\OtherTok{\textless{}{-}}\NormalTok{ (newWeight }\SpecialCharTok{\textgreater{}}\NormalTok{ pred[,}\DecValTok{2}\NormalTok{])}
\NormalTok{smallerthanupper }\OtherTok{\textless{}{-}}\NormalTok{ (newWeight }\SpecialCharTok{\textless{}}\NormalTok{ pred[,}\DecValTok{3}\NormalTok{])}
\NormalTok{inInterval }\OtherTok{\textless{}{-}} \FunctionTok{which}\NormalTok{(largerthanlower }\SpecialCharTok{\&}\NormalTok{ smallerthanupper)}
\FunctionTok{print}\NormalTok{(}\FunctionTok{length}\NormalTok{(inInterval))}
\end{Highlighting}
\end{Shaded}

\begin{verbatim}
## [1] 381
\end{verbatim}

Proportion of observations within prediction intervals:

\begin{Shaded}
\begin{Highlighting}[]
\FunctionTok{length}\NormalTok{(inInterval)}\SpecialCharTok{/}\FunctionTok{length}\NormalTok{(newWeight)}
\end{Highlighting}
\end{Shaded}

\begin{verbatim}
## [1] 0.9844961
\end{verbatim}

What should we expect?
\end{frame}

\end{document}
